% =============================================================================
% Bayesian Trust Update Model: Formal Specification
% =============================================================================
% Trust Calibration Trajectories in Human-AI Interaction
% Paper 4: Theoretical Framework
%
% Author: Hosung You
% Date: 2026-03-12
% =============================================================================

\documentclass[12pt]{article}
\usepackage{amsmath, amssymb, amsthm}
\usepackage{geometry}
\usepackage{booktabs}
\usepackage{graphicx}
\usepackage{natbib}

\geometry{margin=1in}

\newtheorem{definition}{Definition}
\newtheorem{proposition}{Proposition}
\newtheorem{corollary}{Corollary}

\title{Bayesian Trust Update Model:\\
  A Formal Specification of Trust Calibration Dynamics\\
  in Human-AI Interaction}
\author{Hosung You\\
  College of Education, Pennsylvania State University}
\date{March 2026}

\begin{document}
\maketitle

% =============================================================================
\section{Definitions}
% =============================================================================

\begin{definition}[Trust Calibration Space]
  The \textbf{Trust Calibration Space} is a three-dimensional space
  $\mathcal{S} = [0,1] \times [0,1] \times \mathbb{Z}^+$ defined by:
  \begin{itemize}
    \item $T_t \in [0,1]$: \textbf{Behavioral trust} at time $t$,
      operationalized as reliance rate $R_b$
      (proportion of AI recommendations followed)
    \item $R_t \in [0,1]$: \textbf{AI reliability} at time $t$,
      operationalized as AI accuracy
      (proportion of correct AI recommendations)
    \item $\tau \in \mathbb{Z}^+$: \textbf{Discrete time index}
      representing sequential interaction episodes or windows
  \end{itemize}
\end{definition}

\begin{definition}[Calibration Gap]
  The \textbf{calibration gap} at time $t$ is defined as:
  \begin{equation}
    G_t = T_t - R_t
  \end{equation}
  where $G_t > 0$ indicates \textit{over-reliance}
  (trust exceeds reliability),
  $G_t < 0$ indicates \textit{under-reliance}
  (trust falls below reliability),
  and $G_t = 0$ indicates \textit{perfect calibration}.
\end{definition}

\begin{definition}[Prediction Error]
  The \textbf{prediction error} at time $t$ is the discrepancy between
  observed AI reliability and current behavioral trust:
  \begin{equation}
    \delta_t = R_t - T_t = -G_t
  \end{equation}
  Positive prediction error ($\delta_t > 0$) indicates AI performs
  better than expected; negative ($\delta_t < 0$) indicates AI performs
  worse than expected.
\end{definition}

% =============================================================================
\section{Core Model: Prediction Error-Driven Trust Update}
% =============================================================================

\subsection{Symmetric Model}

The simplest form assumes a single learning rate governing trust
adaptation to prediction errors:

\begin{equation}
  \boxed{T_{t+1} = T_t + \alpha \cdot \delta_t = T_t + \alpha(R_t - T_t)}
  \label{eq:symmetric}
\end{equation}

where $\alpha \in [0,1]$ is the \textbf{learning rate} parameter
controlling the speed of trust adaptation.

This is formally equivalent to the Rescorla-Wagner learning rule
\citep{rescorla1972theory} and exponential smoothing in time series
analysis. When $R_t = R$ (constant reliability):

\begin{equation}
  T_t = R + (T_0 - R)(1-\alpha)^t
\end{equation}

\begin{proposition}[Convergence under Constant Reliability]
  When $R_t = R$ for all $t$ and $0 < \alpha \leq 1$:
  \begin{enumerate}
    \item $T_t \to R$ as $t \to \infty$ (convergence to true reliability)
    \item $G_t = (1-\alpha)^t \cdot G_0$ (exponential decay of calibration gap)
    \item Half-life of gap: $\tau_{1/2} = -\ln 2 / \ln(1-\alpha)$
  \end{enumerate}
\end{proposition}

\begin{proof}
  From Equation~\ref{eq:symmetric} with constant $R$:
  $G_{t+1} = T_{t+1} - R = T_t + \alpha(R - T_t) - R = (1-\alpha)(T_t - R) = (1-\alpha)G_t$.
  By induction, $G_t = (1-\alpha)^t G_0$.
  Since $|1-\alpha| < 1$, $G_t \to 0$.
  Setting $(1-\alpha)^{\tau_{1/2}} = 1/2$ yields
  $\tau_{1/2} = -\ln 2 / \ln(1-\alpha)$.
\end{proof}

\subsection{Asymmetric Model (Dual Learning Rate)}

Empirical evidence from automation trust research demonstrates that
trust decreases more rapidly after failures than it increases after
successes \citep{lee1992trust, dzindolet2003role}. We capture this
asymmetry by introducing direction-dependent learning rates:

\begin{equation}
  \boxed{T_{t+1} = T_t +
    \begin{cases}
      \alpha^+ \cdot \delta_t & \text{if } \delta_t \geq 0
        \quad \text{(AI better than expected)} \\
      \alpha^- \cdot \delta_t & \text{if } \delta_t < 0
        \quad \text{(AI worse than expected)}
    \end{cases}}
  \label{eq:asymmetric}
\end{equation}

where:
\begin{itemize}
  \item $\alpha^+ \in [0,1]$: \textbf{Upward learning rate}
    --- speed of trust increase when AI exceeds expectations
  \item $\alpha^- \in [0,1]$: \textbf{Downward learning rate}
    --- speed of trust decrease when AI falls below expectations
\end{itemize}

The \textbf{trust asymmetry ratio} is defined as:
\begin{equation}
  \rho = \frac{\alpha^+}{\alpha^-}
\end{equation}

Following \citet{lee1992trust}, the canonical expectation is $\rho < 1$
(trust is harder to build than to destroy).

% =============================================================================
\section{Derivation of Five Trajectory Patterns}
% =============================================================================

The five trajectory patterns emerge as distinct regions in the
$(\alpha^+, \alpha^-)$ parameter space, conditioned on the
environmental reliability profile $\{R_t\}$.

\subsection{Pattern 1: Convergent ($\alpha^+, \alpha^- > 0$; both sufficient)}

\begin{proposition}[Convergent Trajectory]
  When $\alpha^+ > 0$ and $\alpha^- > 0$, with both learning rates
  sufficiently large relative to the rate of change of $R_t$:
  \begin{equation}
    |G_t| \to 0 \quad \text{as } t \to \infty
  \end{equation}
  regardless of the initial gap $G_0$ or the direction of
  reliability change.
\end{proposition}

\textbf{Parameter region:} $\alpha^+ \geq \alpha_{min}$,
$\alpha^- \geq \alpha_{min}$, where $\alpha_{min}$ depends on the
rate of environmental change.

\textbf{Subtypes:}
\begin{itemize}
  \item \textit{Rapid Convergent}: $\alpha^+, \alpha^- > 0.5$
    ($\tau_{1/2} < 2$ windows)
  \item \textit{Gradual Convergent}: $0.1 < \alpha^+, \alpha^- < 0.5$
    ($\tau_{1/2} > 2$ windows)
\end{itemize}

\textbf{Empirical correspondence:}
Chess Class 1 (Adaptive Reducer, $N=20$, C1) and
Class 3 (Gradual Adaptor, $N=25$, C2).

\textbf{Theoretical basis:}
\citet{lee2004trust} --- trust as adaptive, experience-driven
calibration process.

% ---

\subsection{Pattern 2: Stagnant ($\alpha^+ \approx \alpha^- \approx 0$)}

\begin{proposition}[Stagnant Trajectory]
  When $\alpha^+ \approx 0$ and $\alpha^- \approx 0$:
  \begin{equation}
    T_{t+1} \approx T_t \quad \forall t
  \end{equation}
  \begin{equation}
    G_t \approx T_0 - R_t
  \end{equation}
  The calibration gap is determined entirely by the initial trust
  level $T_0$ and the current reliability $R_t$, with no adaptive
  updating.
\end{proposition}

\textbf{Parameter region:} $\alpha^+, \alpha^- < \epsilon$
for small $\epsilon > 0$.

\textbf{Mechanism:} Cognitive disengagement from the trust
update process. The individual does not attend to or process
AI performance feedback, resulting in behavioral inertia.

\textbf{Boundary condition:} This pattern is distinguishable from
rational non-reliance (where $T_0$ is appropriately low given $R$)
only when $G_t \neq 0$ persists despite observable opportunities
for calibration.

\textbf{Empirical correspondence:}
EdNet Class 4 ($N=451$, $R_b \approx 0.06$--$0.11$, $\Delta G \approx +0.03$).

\textbf{Theoretical basis:}
\citet{parasuraman1997humans} --- automation bias as persistent
behavioral pattern; \citet{goddard2012automation} --- trust inertia.

% ---

\subsection{Pattern 3: Catastrophic ($\alpha^- \approx 0$; $R$ drops sharply)}

\begin{proposition}[Catastrophic Trajectory]
  When $\alpha^- \approx 0$ and a trust violation event occurs
  at time $t^*$ such that $R_{t^*} \ll R_{t^*-1}$:
  \begin{equation}
    T_{t} \approx T_{t^*-1} \quad \forall t > t^*
  \end{equation}
  \begin{equation}
    G_t \approx T_{t^*-1} - R_{t} \gg 0 \quad \forall t > t^*
  \end{equation}
  The calibration gap jumps abruptly to a large positive value
  (over-reliance) and persists because the downward learning
  rate is near zero.
\end{proposition}

\textbf{Parameter region:} $\alpha^- < \epsilon$, $\alpha^+$
unconstrained.

\textbf{Environmental condition required:}
Trust violation event --- a discrete, significant drop in $R_t$.
Without such an event, this pattern is observationally
indistinguishable from Stagnant.

\textbf{Distinction from Stagnant:}
\begin{itemize}
  \item Stagnant: $\alpha^+ \approx \alpha^- \approx 0$
    (bilateral learning failure; no environmental trigger needed)
  \item Catastrophic: $\alpha^- \approx 0$, $\alpha^+$ may be $> 0$
    (unilateral downward learning failure; requires $R$ drop to manifest)
\end{itemize}

\textbf{Empirical correspondence:}
Chess Class 5 ($N=30$, 60\% of C1):
$R_b$ drops only from 0.700 to 0.620 ($\Delta = -0.080$)
despite AI accuracy dropping from 1.000 to 0.200.
$G$ jumps from $-0.300$ to $+0.420$ ($\Delta = +0.720$).

\textbf{Theoretical basis:}
\citet{lee1992trust} --- trust violation and asymmetric recovery;
\citet{parasuraman1993performance} --- automation complacency;
anchoring effect \citep{tversky1974judgment}.

% ---

\subsection{Pattern 4: Oscillating (nonlinear $\alpha$ or $\alpha > 1$)}

\begin{proposition}[Oscillating Trajectory]
  When either (a) $\alpha > 1$ causing overshooting, or
  (b) $\alpha$ is threshold-dependent:
  \begin{equation}
    \alpha_t =
    \begin{cases}
      \alpha_H & \text{if } |G_t| > \theta \\
      \alpha_L & \text{if } |G_t| \leq \theta
    \end{cases}
    \quad \text{with } \alpha_H \gg \alpha_L
  \end{equation}
  the trust trajectory exhibits non-monotonic fluctuation
  around the reliability level, failing to achieve stable
  convergence.
\end{proposition}

\textbf{Mechanism A --- Overshooting:}
When $\alpha > 1$, the update $T_{t+1} = T_t + \alpha(R_t - T_t)$
pushes $T$ past $R$, reversing the sign of $\delta$ and creating
a damped (or sustained) oscillation.

\textbf{Mechanism B --- Threshold switching:}
When large gaps trigger strong corrections but small gaps
are ignored, a limit cycle emerges:
$|G|$ large $\to$ strong correction $\to$ $|G|$ small $\to$
weak maintenance $\to$ $|G|$ grows $\to$ cycle repeats.

\textbf{Note on evidence strength:}
This pattern has the weakest empirical support (GRADE: Low).
With only 6 windows in Chess and 10 in EdNet, distinguishing
true oscillation from measurement noise is difficult.
This pattern should be presented as a theoretical possibility
with preliminary evidence, pending longer time series data.

\textbf{Empirical correspondence:}
Chess Class 2 ($N=14$): non-monotonic $R_b$ trajectory,
$\text{gap\_sd} = 0.421$ (highest variability).

\textbf{Theoretical basis:}
\citet{gao2006extending} --- oscillatory trust recovery;
\citet{muir1996trust} --- repeated failure-recovery cycles.

% ---

\subsection{Pattern 5: AI Benefit Emergence ($\alpha^+ \ll \alpha^-$; $R$ rises)}

\begin{proposition}[AI Benefit Emergence Trajectory]
  When $\alpha^+ \ll \alpha^-$ (trust asymmetry ratio $\rho \ll 1$)
  and AI reliability gradually increases ($R_t \uparrow$ over time):
  \begin{equation}
    \frac{dT}{dt} \approx \alpha^+ \cdot (R_t - T_t) \ll \frac{dR}{dt}
  \end{equation}
  \begin{equation}
    G_t = T_t - R_t < 0 \quad \text{and} \quad |G_t| \uparrow
    \text{ over time}
  \end{equation}
  The under-reliance gap progressively widens as AI improves
  faster than trust adapts.
\end{proposition}

\textbf{Parameter region:} $\alpha^+ < \epsilon$, $\alpha^- > 0$.

\textbf{Environmental condition required:}
Gradual or sudden improvement in $R_t$.
Without reliability improvement, this pattern cannot manifest.

\textbf{Mirror symmetry with Catastrophic:}
\begin{equation}
  \text{Catastrophic:} \quad \alpha^- \approx 0, \; R \downarrow
    \implies G \gg 0 \; \text{(over-reliance)}
\end{equation}
\begin{equation}
  \text{ABE:} \quad \alpha^+ \approx 0, \; R \uparrow
    \implies G \ll 0 \; \text{(under-reliance)}
\end{equation}

These two patterns are \textbf{mathematical mirror images} in the
$(\alpha^+, \alpha^-)$ parameter space, representing unilateral
non-adaptation in opposite directions.

\textbf{Discovery status:}
This pattern was \textit{inductively discovered} during exploratory
GMM analysis of EdNet data. It was not among the four
\textit{a priori} patterns derived from the literature.
Post-hoc theoretical interpretation links it to three established
mechanisms:
\begin{enumerate}
  \item \textbf{Disuse} \citep{parasuraman1997humans}:
    failure to utilize an automated aid despite its reliability
  \item \textbf{Trust asymmetry} \citep{lee1992trust}:
    trust builds more slowly than it decays ($\rho < 1$),
    here in the extreme case ($\rho \to 0$)
  \item \textbf{Anchoring} \citep{tversky1974judgment}:
    initial low trust serves as an anchor resistant to
    upward adjustment
\end{enumerate}

\textbf{Empirical correspondence:}
EdNet S2 Class 4 ($N=256$, 5.7\%): $P_{\text{adaptive}}$
$0.275 \to 0.852$, $R_b$ $0.020 \to 0.171$, $G \to -0.601$.
Chess Class 4 ($N=9$, C2): $R_b = 0.200$, $G = -0.600$.

% =============================================================================
\section{Parameter Space Map}
% =============================================================================

The five patterns occupy distinct regions of the
$(\alpha^+, \alpha^-)$ parameter space (Figure~1):

\begin{center}
\begin{tabular}{lccl}
  \toprule
  \textbf{Pattern} & $\alpha^+$ & $\alpha^-$
    & \textbf{Environmental condition} \\
  \midrule
  Convergent & $> \alpha_{\min}$ & $> \alpha_{\min}$
    & Any (robust to $R$ changes) \\
  Stagnant & $\approx 0$ & $\approx 0$
    & Any (no response to $R$) \\
  Catastrophic & unconstrained & $\approx 0$
    & $R \downarrow$ (trust violation) \\
  ABE & $\approx 0$ & unconstrained
    & $R \uparrow$ (trust repair / AI improvement) \\
  Oscillating & $> 1$ or nonlinear & $> 1$ or nonlinear
    & $R$ variable \\
  \bottomrule
\end{tabular}
\end{center}

\begin{proposition}[Completeness]
  The five patterns partition the behaviorally relevant
  $(\alpha^+, \alpha^-)$ space:
  \begin{enumerate}
    \item Both learning rates active $\to$ Convergent
    \item Both learning rates inactive $\to$ Stagnant
    \item Only upward inactive ($\alpha^+ \approx 0$) $\to$ ABE
      (when $R \uparrow$)
    \item Only downward inactive ($\alpha^- \approx 0$) $\to$ Catastrophic
      (when $R \downarrow$)
    \item Nonlinear / overshooting dynamics $\to$ Oscillating
  \end{enumerate}
\end{proposition}

% =============================================================================
\section{Testable Predictions}
% =============================================================================

The model generates the following testable predictions for
confirmatory studies:

\begin{enumerate}
  \item \textbf{Individual consistency:}
    An individual's $(\alpha^+, \alpha^-)$ values should be
    relatively stable across tasks within a session, predicting
    consistent pattern membership across different AI reliability
    conditions.

  \item \textbf{Asymmetry prevalence:}
    In the general population, $\alpha^- > \alpha^+$ (i.e., $\rho < 1$)
    should be the modal configuration, consistent with
    \citet{lee1992trust}.

  \item \textbf{Catastrophic--ABE mirror:}
    Individuals classified as Catastrophic under trust violation
    ($R \downarrow$) should show elevated ABE risk under trust
    repair ($R \uparrow$), because both patterns share the feature
    of unilateral non-adaptation (one learning rate near zero).

  \item \textbf{Domain expertise moderation:}
    Higher domain expertise should increase both $\alpha^+$ and
    $\alpha^-$ (more accurate performance monitoring enables
    faster trust updating), shifting individuals toward the
    Convergent region.

  \item \textbf{Intervention targets:}
    Calibration interventions (e.g., AI confidence displays,
    performance feedback) should primarily affect the deficient
    learning rate: $\alpha^+$ for ABE-prone individuals,
    $\alpha^-$ for Catastrophic-prone individuals.
\end{enumerate}

% =============================================================================
\section{Scope and Limitations}
% =============================================================================

\begin{enumerate}
  \item \textbf{What $\alpha$ does not explain:}
    The model describes \textit{how} trust updates occur but not
    \textit{why} individuals differ in $\alpha^+$ and $\alpha^-$.
    Determinants of learning rates (personality traits, prior
    experience, cognitive style, domain expertise) are outside
    the current model's scope.

  \item \textbf{Exogenous reliability:}
    The model treats $R_t$ as exogenously given. In interactive
    systems where human behavior affects AI performance (e.g.,
    through data quality), a coupled dynamical system would be
    required.

  \item \textbf{Illustrative status in Paper 4:}
    In the current paper, this model serves as a theoretical
    framework for organizing and interpreting exploratory findings.
    Formal parameter estimation is deferred to subsequent
    experimental studies where trial-level data permits
    individual $\alpha$ fitting.

  \item \textbf{Linearity assumption:}
    The model assumes linear prediction error integration.
    Nonlinear variants (e.g., sigmoid trust functions, Bayesian
    updating with beta priors) may better capture boundary effects
    and extreme reliability regimes.
\end{enumerate}

% =============================================================================
\section*{References}
% =============================================================================

\begingroup
\renewcommand{\section}[2]{}
\begin{thebibliography}{99}

\bibitem[Dzindolet et al.(2003)]{dzindolet2003role}
Dzindolet, M.~T., Peterson, S.~A., Pomranky, R.~A., Pierce, L.~G.,
\& Beck, H.~P. (2003).
The role of trust in automation reliance.
\textit{International Journal of Human-Computer Studies}, 58(6), 697--718.

\bibitem[Gao \& Lee(2006)]{gao2006extending}
Gao, J., \& Lee, J.~D. (2006).
Extending the decision field theory to model operators'
reliance on automation in supervisory control situations.
\textit{IEEE Transactions on Systems, Man, and Cybernetics---Part A},
36(5), 943--959.

\bibitem[Goddard et al.(2012)]{goddard2012automation}
Goddard, K., Roudsari, A., \& Wyatt, J.~C. (2012).
Automation bias: A systematic review of frequency, effect mediators,
and mitigators.
\textit{Journal of the American Medical Informatics Association},
19(1), 121--127.

\bibitem[Lee \& Moray(1992)]{lee1992trust}
Lee, J.~D., \& Moray, N. (1992).
Trust, control strategies and allocation of function in
human-machine systems.
\textit{Ergonomics}, 35(10), 1243--1270.

\bibitem[Lee \& See(2004)]{lee2004trust}
Lee, J.~D., \& See, K.~A. (2004).
Trust in automation: Designing for appropriate reliance.
\textit{Human Factors}, 46(1), 50--80.

\bibitem[Muir \& Moray(1996)]{muir1996trust}
Muir, B.~M., \& Moray, N. (1996).
Trust in automation. Part II. Experimental studies of trust and
human intervention in a process control simulation.
\textit{Ergonomics}, 39(3), 429--460.

\bibitem[Parasuraman \& Riley(1997)]{parasuraman1997humans}
Parasuraman, R., \& Riley, V. (1997).
Humans and automation: Use, misuse, disuse, abuse.
\textit{Human Factors}, 39(2), 230--253.

\bibitem[Parasuraman et al.(1993)]{parasuraman1993performance}
Parasuraman, R., Molloy, R., \& Singh, I.~L. (1993).
Performance consequences of automation-induced ``complacency.''
\textit{The International Journal of Aviation Psychology}, 3(1), 1--23.

\bibitem[Rescorla \& Wagner(1972)]{rescorla1972theory}
Rescorla, R.~A., \& Wagner, A.~R. (1972).
A theory of Pavlovian conditioning: Variations in the effectiveness
of reinforcement and nonreinforcement.
In A.~H. Black \& W.~F. Prokasy (Eds.),
\textit{Classical conditioning II: Current research and theory}
(pp. 64--99). Appleton-Century-Crofts.

\bibitem[Tversky \& Kahneman(1974)]{tversky1974judgment}
Tversky, A., \& Kahneman, D. (1974).
Judgment under uncertainty: Heuristics and biases.
\textit{Science}, 185(4157), 1124--1131.

\end{thebibliography}
\endgroup

\end{document}
